\section{本章小结}
\label{sec:graph-laplacian-summary}

组合拉普拉斯把沿边取差、乘权和返回顶点三个步骤合成为$\bm L=\bm B\bm C\bm B^\top$。逐点来看，它输出的是每个顶点与邻居的带权偏离之和。若要衡量整张图上读数的变化而不让正负相消，就逐边取差、平方、乘权再求和。这个沿边变化能量恰好等于拉普拉斯的二次型$\bm x^\top\bm L\bm x$。

能量为零要求每条正权边的两端读数相等。沿正权路径传递这个等值关系，每个连通分量上的读数各为常值，分量之间可以独立选择。分量的个数就是零特征值的重数。

按边集展开约化拉普拉斯的行列式时，只有构成生成树的边组留下非零的关联子式。因此，删去任一个顶点对应的行和列后所得的主余子式都给出同一个树权总和。对正权支撑连通的图，也可以用$n-1$个非零特征值的乘积除以顶点数$n$来计算它。

每个特征值还描述了对应单位特征向量的变化尺度——长度为一的读数方向上的沿边能量。加边使每个方向的单位能量不减，进而使按大小排序的各特征值逐个不减。

三种矩阵使用不同的内积和坐标。组合拉普拉斯$\bm L=\bm D-\bm W$在原顶点读数上计算带权差的总和。随机游走拉普拉斯$\bm L_{\mathrm{rw}}=\bm I-\bm P$用同样的原顶点读数，输出的是本点读数与邻居加权平均的差。对称归一化拉普拉斯$\mathcal L=\bm D^{-1/2}\bm L\bm D^{-1/2}$在缩放坐标$\bm y=\bm D^{1/2}\bm x$上计算，使得度加权能量比变成普通 Rayleigh 商。三种形式的零方向都来自各分量上的常值读数——在原坐标中是分量指示向量$\bm1_{V_a}$，在缩放坐标中变为$\bm D^{1/2}\bm1_{V_a}$。三种矩阵都对孤点留出零块。

后续各章使用的矩阵形式各不相同。第~\ref{chap:graph-electrical-networks}章的电网络从注入电流出发求电势，方程的系数是组合拉普拉斯。第~\ref{chap:spectral-partition}章要比较连接量不同的顶点集合两侧的割与体积，归一化形式提供了按体积计量的尺度。第~\ref{chap:graph-walks}章用转移矩阵$\bm P$描述逐步选择邻居的演化，从邻居平均出发的收敛分析建立在归一化谱之上。
